\section{Broader Impacts}
We adopt and improve a framework of generative agents introduced in \citet{park2023generative} for the generation of long-term conversations. Consequently, the ethical concerns of generative agents outlined by \citet{park2023generative} apply to our work as well, especially since the goal of our framework is to make the conversations as realistic as possible. 


Specifically, conversational agents that can pose as human beings with a realistic life, as enabled by the temporal event graphs in our framework, pose the risk that users may form parasocial relationships with such agents that may affect their lives adversely. We recommend that any practical deployment of the generative frameworks mentioned in our work be always prefaced with a disclaimer about the source of the dialogs.


Second, the use of multimodal LLMs \cite{zheng2023minigpt} to generate images conditioned on dialog can lead to the propagation of misinformation and social biases, especially if the conversational agent can be coerced into parroting false information or dangerous opinions.

Third, it is tempting to use generative agents to substitute real humans for a process, especially when there are significant challenges in working with humans for a particular goal e.g., collecting real-world interactions between humans over a year or more. Care must be taken to ensure that such substitutes are not made in studies whose outcomes may be used to make real-world decisions with tangible impacts on humans. Our work is merely a study of model comprehension in very long-term conversations. We do not make any recommendations for real-world policies based on this study and advise potential users of our framework to avoid making such recommendations as well.
